\documentclass[10pt,a4paper]{amsart}
\usepackage{amsmath,amssymb,amsthm,graphicx}
\usepackage{verbatim}
\newenvironment{reference}{\comment}{\endcomment}
%\newenvironment{reference}{}{}
\newenvironment{slogan}{\comment}{\endcomment}
\newenvironment{history}{\comment}{\endcomment}
% For commutative diagrams we use Xy-pic
\usepackage[all]{xy}

% We use 2cell for 2-commutative diagrams.
\xyoption{2cell}
% Theorem environments.
%
\theoremstyle{plain}
\newtheorem{theorem}[subsection]{Theorem}
\newtheorem{proposition}[subsection]{Proposition}
\newtheorem{lemma}[subsection]{Lemma}

\theoremstyle{definition}
\newtheorem{definition}[subsection]{Definition}
\newtheorem{example}[subsection]{Example}
\newtheorem{exercise}[subsection]{Exercise}
\newtheorem{situation}[subsection]{Situation}

\theoremstyle{remark}
\newtheorem{remark}[subsection]{Remark}
\newtheorem{remarks}[subsection]{Remarks}

\numberwithin{equation}{subsection}

% Macros
%
\def\lim{\mathop{\mathrm{lim}}\nolimits}
\def\colim{\mathop{\mathrm{colim}}\nolimits}
\def\Spec{\mathop{\mathrm{Spec}}}
\def\Hom{\mathop{\mathrm{Hom}}\nolimits}
\def\Ext{\mathop{\mathrm{Ext}}\nolimits}
\def\SheafHom{\mathop{\mathcal{H}\!\mathit{om}}\nolimits}
\def\SheafExt{\mathop{\mathcal{E}\!\mathit{xt}}\nolimits}
\def\Sch{\mathit{Sch}}
\def\Mor{\mathop{\mathrm{Mor}}\nolimits}
\def\Ob{\mathop{\mathrm{Ob}}\nolimits}
\def\Sh{\mathop{\mathit{Sh}}\nolimits}
\def\NL{\mathop{N\!L}\nolimits}
\def\CH{\mathop{\mathrm{CH}}\nolimits}
\def\proetale{{pro\text{-}\acute{e}tale}}
\def\etale{{\acute{e}tale}}
\def\QCoh{\mathit{QCoh}}
\def\Ker{\mathop{\mathrm{Ker}}}
\def\Im{\mathop{\mathrm{Im}}}
\def\Coker{\mathop{\mathrm{Coker}}}
\def\Coim{\mathop{\mathrm{Coim}}}

% Boxtimes
%
\DeclareMathSymbol{\boxtimes}{\mathbin}{AMSa}{"02}

%
% Macros for moduli stacks/spaces
%
\def\QCohstack{\mathcal{QC}\!\mathit{oh}}
\def\Cohstack{\mathcal{C}\!\mathit{oh}}
\def\Spacesstack{\mathcal{S}\!\mathit{paces}}
\def\Quotfunctor{\mathrm{Quot}}
\def\Hilbfunctor{\mathrm{Hilb}}
\def\Curvesstack{\mathcal{C}\!\mathit{urves}}
\def\Polarizedstack{\mathcal{P}\!\mathit{olarized}}
\def\Complexesstack{\mathcal{C}\!\mathit{omplexes}}
% \Pic is the operator that assigns to X its picard group, usage \Pic(X)
% \Picardstack_{X/B} denotes the Picard stack of X over B
% \Picardfunctor_{X/B} denotes the Picard functor of X over B
\def\Pic{\mathop{\mathrm{Pic}}\nolimits}
\def\Picardstack{\mathcal{P}\!\mathit{ic}}
\def\Picardfunctor{\mathrm{Pic}}
\def\Deformationcategory{\mathcal{D}\!\mathit{ef}}
\usepackage{cleveref}
\crefname{originallemma}{original lemma}{original lemmas}
\Crefname{originallemma}{Original lemma}{Original lemmas}
\makeatletter
\newlabel{lemma-finite-ext-into-bounded-injective}{{lemma-finite-ext-into-bounded-injective}{source}{}{sourcepointer-1}{}}
\newlabel{lemma-finite-ext-into-bounded-injective@cref}{{[originallemma][0][]lemma-finite-ext-into-bounded-injective}{[1][0][]source}{}{}{}}
\makeatother
\crefname{originalreference}{original reference}{original references}
\Crefname{originalreference}{Original reference}{Original references}
\makeatletter
\newlabel{more-algebra-lemma-internal-hom-evaluate-isomorphism-technical}{{more-algebra-lemma-internal-hom-evaluate-isomorphism-technical}{source}{}{sourcepointer-2}{}}
\newlabel{more-algebra-lemma-internal-hom-evaluate-isomorphism-technical@cref}{{[originalreference][0][]more-algebra-lemma-internal-hom-evaluate-isomorphism-technical}{[1][0][]source}{}{}{}}
\makeatother
\title{The dualizing functor is an anti-autoequivalence of the bounded coherent category with square the identity}
\author{The Stacks Project Authors}
\date{stacks-project:dualizing.tex:local-snapshot-20261012; GNU FDL 1.2 or later}
\begin{document}
\maketitle
\noindent\textit{Editorial note.} All delivered mathematical text is the author's own, copied from the pinned Stacks Project chapter source sources/stacks-dualizing/pinned-original-source.tex; no statement or proof step was rewritten, reordered or removed.\par\medskip
\noindent\textit{Editorial note.} Preamble provenance: the Stacks Project keeps one single shared preamble for every chapter; this card takes it directly from the staged sources/stacks-dualizing/preamble.tex as three verbatim spans (lines 8-12, 14-18 and 154-217), so no derived file was written and nothing was added or redefined. The wrapper class and the base packages are supplied by the assembler.\par\medskip
\noindent\textit{Editorial note.} Bibliography provenance: neither the counted statement nor its proof contains a citation (this lane rescans the whole span before writing the card), so the card carries no bibliography and no entry was reconstructed.\par\medskip
\noindent\textit{Editorial note.} Selection provenance: this card was opened under the file-based control contract (CONTROL.md), with the mechanical pre-card gate run on its statement and proof interval before the card was written.\par\medskip
\noindent\textit{Allowed standard prerequisites (editorial).}\begin{itemize}
\item The statement that the internal hom into a dualizing complex preserves the bounded coherent derived category. Invoked at the start of the counted proof (source lines 2832-2834). Shows the functor lands where claimed.
\item The evaluation isomorphism identifying an object with the tensor of the double dual of the dualizing complex. Invoked inside the counted proof (source lines 2835-2840) as a cross-chapter pointer. Produces the canonical comparison map to the double dual.
\end{itemize}\medskip
\ifdefined\hypertarget\hypertarget{sourcepointer-1}{}\fi
\noindent\textit{Source reference.} Original lemma lemma-finite-ext-into-bounded-injective: sources/stacks-dualizing/pinned-original-source.tex, source lines 2791--2791. Navigation pointer only: the counted text refers to the Stacks item lemma-finite-ext-into-bounded-injective, which is an existing labeled result of this same chapter and is not part of this card. The pointer carries the target's own label as its number and its definition line as its span; its content is not reproduced here.\par\medskip
\ifdefined\hypertarget\hypertarget{sourcepointer-2}{}\fi
\noindent\textit{Source reference.} Original reference more-algebra-lemma-internal-hom-evaluate-isomorphism-technical: sources/stacks-dualizing/pinned-original-source.tex, source lines 2836--2836. Navigation pointer only: the counted text refers to the Stacks item more-algebra-lemma-internal-hom-evaluate-isomorphism-technical, which belongs to another chapter of the project and is not part of this card. The pointer carries the target's own label as its number because this chapter's staged snapshot does not contain the other chapter's tag table; its content is not reproduced here.\par\medskip
\begin{lemma}
\label{lemma-dualizing}
Let $A$ be a Noetherian ring. If $\omega_A^\bullet$ is a dualizing
complex, then the functor
$$
D : K \longmapsto R\Hom_A(K, \omega_A^\bullet)
$$
is an anti-equivalence $D_{\textit{Coh}}(A) \to D_{\textit{Coh}}(A)$
which exchanges $D^+_{\textit{Coh}}(A)$ and $D^-_{\textit{Coh}}(A)$
and induces an anti-equivalence
$D^b_{\textit{Coh}}(A) \to D^b_{\textit{Coh}}(A)$.
Moreover $D \circ D$ is isomorphic to the identity functor.
\end{lemma}

\begin{proof}
Let $K$ be an object of $D_{\textit{Coh}}(A)$. From
Lemma \ref{lemma-finite-ext-into-bounded-injective}
we see $R\Hom_A(K, \omega_A^\bullet)$ is an object of $D_{\textit{Coh}}(A)$.
By More on Algebra, Lemma
\ref{more-algebra-lemma-internal-hom-evaluate-isomorphism-technical}
and the assumptions on the dualizing complex
we obtain a canonical isomorphism
$$
K = R\Hom_A(\omega_A^\bullet, \omega_A^\bullet) \otimes_A^\mathbf{L} K
\longrightarrow
R\Hom_A(R\Hom_A(K, \omega_A^\bullet), \omega_A^\bullet)
$$
Thus our functor has a quasi-inverse and the proof is complete.
\end{proof}

\end{document}
